\begin{definition}
	\textit{Евклидовым пространством} над полем $\mathbb{K}$( где $\mathbb{K} = \R$ или $\mathbb{K} = \C$), называется линейное пространство $E$ над \ $\mathbb{K}$ \ с функцией $(\cdot, \cdot) : E^2 \to \mathbb{K}$, обладающей следующими свойствами:
	\begin{enumerate}
		\item $\forall x \in E: (x, x) \geqslant 0$, причем $(x, x) = 0 \lra x = 0$
		\item $\forall x, y \in E: (x, y) = \overline{(y, x)}$
		\item $\forall x, y, z \in E: \forall \alpha, \beta \in \mathbb{K} (\alpha x + \beta y, z) = \alpha(x, z) + \beta(y, z)$
	\end{enumerate}

	Функция $(\cdot, \cdot)$ называется \textit{скалярным произведением} на пространстве $E$.
\end{definition}

\begin{example}
	Легко видеть, что любое евклидово пространство $E$ является линейным нормированным с нормой, заданной для произвольного $x \in E$ как $\|x\| := \sqrt{(x, x)}$. Некоторые из приведенных ранее в качестве примеров метрических пространств являются не только линейными нормированными, но и евклидовыми, и норма в них порождается скалярным произведением:
	\begin{enumerate}
		\item Скалярное произведение на $\R^n_2$ задается для произвольных $x, y \in \R^n$ следующим образом:
  		\[(x, y) := \sum_{k = 1}^n x_ky_k\]

		\item Скалярное произведение на $l_2$ задается для произвольных $x, y \in l_2$ следующим образом:
		\[(x, y) := \sum_{n = 1}^\infty x_ny_n\]

		\item Скалярное произведение на $L_2[a, b]$ задается для произвольных $f, g \in L_2[a, b]$ следующим образом:
		\[(f, g) := \int_a^b f(x)g(x)dx\]
	\end{enumerate}
\end{example}

\begin{theorem}[Характеристическое свойство Евклидовых пространств (б/д)]
	Пусть $E$ "--- линейное нормированное пространство. Тогда норма в $E$ порождается скалярным произведением $\lra$ для любых $x, y \in E$ выполнено равенство параллелограмма:
	\[\|x+y\|^2 + \|x-y\|^2 = 2\|x\|^2 + 2\|y\|^2\]
\end{theorem}

\begin{definition}~ Полное линейное нормированное пространство $E$ называется \textit{банаховым}.
\end{definition}

\begin{definition}~ Полное евклидово пространство $H$ называется \textit{гильбертовым}.
\end{definition}

\begin{corollary}[неравенство Коши---Буняковского]
	Пусть $E$ "--- евклидово пространство. Тогда для любых $x, y \in E$ выполнено следующее неравенство:
	\[|(x, y)| \leqslant \|x\| \cdot \|y\|\]
\end{corollary}

\begin{proof}
$e = \frac{y}{\|y\|}$ - единичный вектор
 $$0 \leqslant \|x - (x,e)e\|^2 = \|x\|^2 - \overline{(x,e)}(x,e) - (x,e)(e,x) + |(x,e)|^2 = \|x\|^2 - |(x,e)|^2$$
 $$|(x,e)| \leqslant \|x\|$$
 $$|(x,e)|\cdot\|y\| = |(x,y)| \leqslant \|x\| \cdot \|y\|$$
\end{proof}

\begin{note}
	Из неравенства Коши-Буняковского, в частности, следует, что скалярное произведение непрерывно на $E$ относительно порождаемой им нормы.
\end{note}